\section*{Bab \ref{ch:g9:immune-defenses} --- Pertahanan Kekebalan}

\begin{solution}{exo:g9:immune-defenses:1}
\omterm{def:g9:immune-defenses:standing}{Sel darah putihnya}, yang dibuat di dalam sumsum \omterm{def:g3:bones-and-muscles:skeleton}{tulangnya}
lalu berpatroli di \omterm{prop:g4:journey-of-food:absorb}{darah} dan jaringannya. Penanggap
pertamanya: para penelan, yang mengalir mengelilingi
penyusupnya lalu mencernanya.
\end{solution}

\begin{solution}{exo:g9:immune-defenses:2}
Kemerahan dan kehangatannya: \omterm{def:g4:heart-and-blood:vessels}{pembuluhnya} melebar, lebih
banyak \omterm{prop:g4:journey-of-food:absorb}{darah} di lokasinya. Pembengkakannya: dindingnya bocor
sehingga cairan dan pasukannya keluar. Denyut nyerinya:
medan perang yang sesak dan penuh isyarat. (Dan nanahnya,
bila datang: daftar korban penelan yang gugur dan penyusup
yang mati.)
\end{solution}

\begin{solution}{exo:g9:immune-defenses:3}
\omterm{prop:g9:immune-defenses:specific}{Antigen} adalah tanda permukaan seorang penyusup --- sebuah
bentuk molekul yang dibaca sebagai asing. \omterm{prop:g9:immune-defenses:specific}{Antibodi} adalah
molekul yang dipaskan lalu dicurahkan ke dalam \omterm{prop:g4:journey-of-food:absorb}{darah} oleh
para ahlinya. Hubungannya: bentuk \omterm{prop:g9:immune-defenses:specific}{antibodinya} cocok dengan
\omterm{prop:g9:immune-defenses:specific}{antigen} itu, lalu memakukan dan menggumpalkan pengangkutnya
bagi para penelannya.
\end{solution}

\begin{solution}{exo:g9:immune-defenses:4}
Karena sebuah \omterm{def:g9:microbes-and-infection:sorted}{virus} memperbanyak diri \emph{di dalam} \omterm{def:g5:first-look-at-cells:cell}{sel}
tubuhnya, di atas permesinan \omterm{def:g5:first-look-at-cells:cell}{sel} itu: tidak ada serangan
atas \omterm{def:g9:microbes-and-infection:sorted}{virus} yang bebas yang menghentikan pencetakannya.
Memusnahkan sel-mesin-cetak yang dibajak adalah satu-satunya
cara menghentikan produksinya.
\end{solution}

\begin{solution}{exo:g9:immune-defenses:5}
Seorang ahli berumur panjang yang memegang perian seorang
penyusup yang sudah dikalahkan. Sebuah pendaratan kedua
bertemu korps yang siap: tanggapan dalam hitungan jam,
dibersihkan sebelum gejalanya --- itulah \omterm{def:g9:immune-defenses:memory}{kekebalan}.
\end{solution}

\begin{solution}{exo:g9:immune-defenses:6}
\omterm{prop:g9:immune-defenses:vaccination}{Vaksinasi} menyodorkan kepada para ahli \omterm{def:g9:immune-defenses:memory}{kekebalannya} sebuah
perian penyusup yang tak merugikan, sehingga ingatannya
tercipta tanpa penyakitnya.
\end{solution}

\begin{solution}{exo:g9:immune-defenses:7}
Karena banyak \omterm{prop:g9:microbes-and-infection:mostly}{patogen} memperbanyak diri lebih buruk dalam
keadaan hangat --- aturan \omterm{def:g6:exploring-our-environment:conditions}{kondisi} pada
\cref{prop:g6:decomposers-and-soil:speed} dibalikkan
terhadap penyusupnya: tubuhnya dengan sengaja berjalan panas
untuk memperlambat pelipatduaan musuhnya selagi tentara
klonanya dibangkitkan. Ongkos, ya; kerusakan, tidak.
\end{solution}

\begin{solution}{exo:g9:immune-defenses:8}
\omterm{prop:g9:immune-defenses:specific}{Antigen} campaknya mantap: ingatan satu kemenangan cocok
dengan tiap perjumpaan berikutnya --- sekali campak, tidak
pernah dua kali. \omterm{def:g9:microbes-and-infection:sorted}{Virus} pileknya bermutasi cepat: galur tiap
musim dinginnya memakai perian yang tidak dipegang
ingatannya, sehingga tiap pilek adalah sebuah pendaratan
pertama --- pilek selamanya, walaupun masing-masing tetap
diingat.
\end{solution}

\begin{solution}{exo:g9:immune-defenses:9}
Keragamannya: para ahli dengan reseptor berbeda yang tak
terhitung banyaknya berdiri siap. Penahannya: persentuhan
\omterm{prop:g9:immune-defenses:specific}{antigennya} --- hanya segelintir yang cocok yang
``terseleksi''. Perbanyakan yang tidak merata: yang
terseleksi memperbanyak diri menjadi tentara klonanya, dan
mewarisi reseptor yang cocok itu. Ketiga kenyataan
seleksinya, dijalankan dalam hitungan hari di dalam satu tubuh.
\end{solution}

\begin{solution}{exo:g9:immune-defenses:10}
Sebuah wabah membutuhkan inang yang dapat tertular: tiap
tetangga yang divaksinasi adalah jalan buntu bagi
penularannya, dan cukup banyak jalan buntu melaparkan
\omterm{def:g6:how-plants-spread:dispersal}{penyebarannya} sebelum ia mencapai semua orang --- itulah
padamnya. Yang bergantung padanya: mereka yang tak dapat
dilatih --- yang terlalu muda, yang terlalu sakit --- yang
berdiri di dalam perisai berkumpulnya atau tidak di mana pun.
\end{solution}

\begin{solution}{exo:g9:immune-defenses:11}
(a) Pendaratan-kedua-yang-diatur: suntikan ulangnya
menyodorkan ulang perian tetanusnya, mengisi ulang sebuah
ingatan yang memudar --- pencegahan yang diatur waktunya
menurut sebuah luka masuk yang baru. (b) Pendaratan pertama,
tahap para ahlinya: \omterm{def:g8:hormonal-communication:glands}{kelenjarnya} adalah baraknya, yang sesak
oleh tentara klona yang sedang dibangkitkan terhadap
penyusup tenggorokannya --- pertahanan yang bekerja, bukan
bencana.
\end{solution}

\begin{solution}{exo:g9:immune-defenses:12}
Virusnya: sebuah teks yang hanya meminjam mesin cetak
manusia --- tanpa waduk \omterm{def:g1:animals-around-us:animal}{hewan}, tanpa tempat lain untuk
bersembunyi. Ingatannya: korps terlatih Jenner, yang
disempurnakan menjadi sebuah vaksin yang mantap dan
perlindungannya seumur hidup. Pengumpulannya: sebuah
kampanye sedunia yang memvaksinasi cincin demi cincin sampai
tidak ada pendaratan yang menemukan inang yang dapat
tertular --- penularannya dilaparkan di mana-mana sekaligus.
Pembasmian adalah kepunahan: spesies-hilang-selamanya pada
\cref{prop:g5:biodiversity:loss}, yang dicapai dengan
sengaja, yang tidak ditangisi seorang pun --- satu-satunya
kepunahan yang diusahakan umat manusia, dan bukti tentang
apa yang dapat dilakukan ingatan yang berkumpul.
\end{solution}

\begin{solution}{pb:g9:immune-defenses:1}
\textbf{1.} Perian para penyusupnya --- \omterm{prop:g9:immune-defenses:specific}{antigen} pada \omterm{def:g6:producing-our-food:fermentation}{mikroba}
yang dilemahkan dan pada serpihan yang tak merugikan.
Program seleksi-dan-ingatannya yang lengkap: ahli yang cocok
terseleksi, tentara klonanya dibangkitkan, \omterm{def:g9:immune-defenses:memory}{sel ingatannya}
mundur menjadi sebuah korps yang terlatih.

\textbf{2.} Membersihkannya melayani pertahanan pertamanya
--- pintunya disapu pada tahap masuknya. Suntikan ulangnya
menyodorkan ulang periannya lalu mengisi ulang ingatannya
sampai kekuatan penuh. Tetanus tidak mengizinkan taruhan
karena pertempuran pertamanya dilawankan dengan sebuah racun
yang melumpuhkan dari luka yang dalam dan tanpa udara ---
jenis pertempuran yang dikalahkan tubuh; pelatihannya harus
mendahuluinya.

\textbf{3.} Hari satu sampai tiga: tanggapan cepatnya
terkejar, virusnya mencetak menembus saluran napas dan
\omterm{def:g1:the-five-senses:sense}{kulitnya}. Demam dan ruam pekan pertamanya: tentara klonanya
dibangkitkan --- ongkos pertempurannya dipakai secara
terlihat. Pekan kedua: \omterm{prop:g9:immune-defenses:specific}{antibodi} dan pembunuhnya membersihkan
mesin cetaknya; \omterm{def:g7:organs-at-work:recovery}{pemulihannya}. Yang mundur: sebuah garnisun
ingatan seumur hidup --- sekali campak, tidak pernah dua kali.

\textbf{4.} Yang divaksinasi: terlatih --- pendaratan apa pun
adalah pendaratan kedua, dibersihkan dalam hitungan jam.
Yang sudah sembuh: kedudukan yang sama, dengan ingatan dari
sebuah pertempuran pertama yang sungguhan. Bayinya: tak
terlatih --- sebuah pendaratan pertama akan menjalankan
program penuh yang berbahaya; karena itulah pengawasan
cemasnya.

\textbf{5.} Bahwa ingatan yang bersesuaian masih berdiri:
kadar \omterm{prop:g9:immune-defenses:specific}{antibodinya} adalah tanda tangan garnisunnya di dalam
\omterm{prop:g4:journey-of-food:absorb}{darah}, yang menyiratkan \omterm{def:g9:immune-defenses:memory}{sel ingatan} yang tersedia dan siap
mencurahkan lebih banyak begitu bersentuhan.

\textbf{6.} Garnisun ingatannya menipis seiring tahun ---
kadarnya memudar di bawah garis pembersihan cepatnya. Sebuah
suntikan ulang adalah penyodoran yang segar: ingatan yang
masih tersisa memperbanyak diri lagi, lalu mengisi ulang
garnisunnya sampai kuat.

\textbf{7.} \omterm{prop:g9:genes-and-dna:explains}{Mutasinya} yang cepat: galur tiap musimnya
mengangkut perian yang dieja ulang, sehingga pelatihan tahun
lalu tidak cocok dengan \omterm{prop:g9:immune-defenses:specific}{antigen} tahun ini --- korpsnya harus
dilatih ulang tiap tahun menurut perian yang sekarang.

\textbf{8.} Karena stafnya adalah penyambung gedungnya ---
setiap hari berada di samping orang sakit dan orang yang
rapuh, justru mereka yang pelatihannya tidak ada atau
melemah. Staf yang divaksinasi adalah jalan buntu di antara
ranjangnya: perlindungan mereka yang berkumpul memerisai
mereka yang paling bergantung di bangsalnya.

\textbf{9.} Ongkos pelatihannya: sebuah \omterm{def:g1:my-body:parts}{lengan} yang pegal
dan sehari menggerutu, dengan ingatan yang diperoleh tanpa
pertempuran. Ongkos pertempuran pertamanya, bagi campak:
sepekan demam tinggi paling baiknya, dengan sebagian nyata
berupa radang paru, dan yang lebih buruk --- dijalankan pada
seorang anak yang sama sekali tanpa ingatan. ``Alami''
menamai jalan yang mahal menuju ingatan yang sama, dikurangi
keselamatannya.

\textbf{10.} Mereka yang divaksinasi di sekelilingnya ---
keluarga, kelasnya, kotanya: perlindungan yang berkumpul,
perisai yang menaungi mereka yang terlalu muda untuk
dilatih. Nama kewajibannya pada proposisinya: berkumpulnya
perlindungan --- perisai yang dipegang di atas mereka yang
tak dapat dilatih.

\textbf{11.} Hari Selasa menjalankan garis Pasteur:
penyodoran yang dilemahkan-atau-dimatikan dengan sengaja,
yang disempurnakan menjadi suntikan ulang. Kontak yang
diawasi dan wabah yang melapar pada hari Rabu menjalankan
garis Jenner: ingatan terlatih yang berdiri di antara sebuah
\omterm{def:g9:microbes-and-infection:sorted}{virus} dan inangnya yang berikutnya.

\textbf{12.} ``Pusatnya memberikan ingatan: terlatih,
disuntik ulang, dan berkumpul --- satu-satunya obat yang
bekerja sebelum musuhnya tiba.''
\end{solution}
